\documentclass[12pt]{article}
\usepackage[utf8]{inputenc}
\usepackage[T1]{fontenc}
\usepackage{amsmath, amssymb, geometry, booktabs}
\usepackage{enumitem}
\geometry{margin=1in}
\setlength{\parindent}{0pt}
\renewcommand{\arraystretch}{1.2}

\title{The Standard Trade Model and External Economies of Scale}
\author{ECON 308 -- International Economic Relations}
\date{Fall 2025}

\begin{document}
\maketitle




\subsection*{Problem 1 – Terms of Trade and Welfare (Standard Trade Model)}

A country produces \textbf{food (F)} and \textbf{manufactures (M)}.  
Its PPF is:
\[
F^2 + M^2 = 400.
\]
The world relative price is \( P_M / P_F = 2. \)

\begin{enumerate}[label=(\alph*)]
\item Determine the country’s production point that maximizes the value of output.  
\item Suppose consumers prefer \( M = 0.5F \). Find the consumption point.  
\item Identify whether the country exports or imports manufactures.  
\item Explain how an increase in \( P_M / P_F \) affects the country’s welfare.  
\end{enumerate}

\vspace{1em}

\subsection*{Problem 2 – Export- and Import-Biased Growth}

Country A exports steel and imports textiles.  
Initially, the world relative price is \( P_S / P_T = 1. \)

\begin{enumerate}[label=(\alph*)]
\item Using a standard trade diagram (PPF + indifference curves), show the effects of \textbf{export-biased growth} in Country A on:
\begin{itemize}
    \item The production point,  
    \item The relative supply of steel, and  
    \item The world relative price.
\end{itemize}
\item Discuss qualitatively whether this growth improves or worsens Country A’s terms of trade.  
\item Repeat for \textbf{import-biased growth} and explain the welfare implications.  
\end{enumerate}



\end{document}
